%% LyX 1.2 created this file.  For more info, see http://www.lyx.org/.
%% Do not edit unless you really know what you are doing.
\documentclass[english,ngerman]{article}
\usepackage[T1]{fontenc}
\usepackage[latin1]{inputenc}
\usepackage{longtable}
\usepackage{amsmath}
\usepackage{a4wide}
\usepackage{babel}
\IfFileExists{url.sty}{\usepackage{url}}
                      {\newcommand{\url}{\texttt}}

\makeatletter

%%%%%%%%%%%%%%%%%%%%%%%%%%%%%% LyX specific LaTeX commands.
\providecommand{\LyX}{L\kern-.1667em\lower.25em\hbox{Y}\kern-.125emX\@}
\newcommand{\noun}[1]{\textsc{#1}}

%%%%%%%%%%%%%%%%%%%%%%%%%%%%%% Textclass specific LaTeX commands.
 \newenvironment{lyxcode}
   {\begin{list}{}{
     \setlength{\rightmargin}{\leftmargin}
     \raggedright
     \setlength{\itemsep}{0pt}
     \setlength{\parsep}{0pt}
     \normalfont\ttfamily}%
    \item[]}
   {\end{list}}

%%%%%%%%%%%%%%%%%%%%%%%%%%%%%% User specified LaTeX commands.

\makeatletter

\makeatother
\begin{document}

\title{Guide - Changes}


\author{Volker Maag, Thomas Willhalm}

\maketitle

\section{Guide - short description}

Guide is a program to visualize large graphs%
\footnote{On an AMD Athlon 1400MHz/256MB RAM it worked with a graph with about
25 000 vertices.%
} efficiently. It has an interactive display to visualize the graph
and it is based on the previous version named GRIP. It can read files in
the gml-format and there are some test graphs like grid and moebius band
implemented. The algorithm uses an intelligent clustering of the vertices
in several layers and a force-based model to arrange the vertices
in a layer in two, three or four dimensions. The force-based model is
based on the method of Kamada-Kawai. Only in the last layer
the method of Fruchtermann-Reingold is used.

You may find the original source code, documentation and user's guide
on \url{http://www.cs.arizona.edu/people/kobourov/GRIP/} 


\section{Extensions}

The following extensions were implemented:

\begin{enumerate}
\item read and write support for the GraphML-format%
\footnote{For an up-do-date documentation of the GraphML-standard look
at \url{http://www.graphdrawing.org/graphml/index.html}%
}
\item support for weighted graphs
\end{enumerate}

\section{Algorithmic Details}


\subsection{Improvements}

While testing it turned out that two-dimensional grids with more than
100 vertices were not drawn as smooth as possible but had some
waves when embedded in three dimensions. This undesirable behaviour
was explored with the intention to eliminate it. However, the algorithm could
not be improved significantly. 


\subsection{Support for weighted graphs}

In the original version edges had only constant weight. By manipulating
this constant the drawing could be zoomed. Now the data structure
for the graph is extended to handle weighted edges. Furthermore, some
weighted test-graphs were implemented. It was neccessary to adjust
the algorithm in following points:

\begin{itemize}
\item Replacement of the simple breadth first search by Djikstra's algorithm
\item Adjustment of some parameters in the calculation of the Fruchtermann-Reingold-/
Kamada-Kawai-vector. The new formulas look like:\newline
Kamada-Kawai%
\footnote{The formula is also used when a vertex is placed for the
  first time but then only the nearest three vertices
in the current layer are considered instead of all vertices.%
}:

\[
disp(u)=\sum _{{\textrm{\tiny all vertices }v\atop \textrm{\tiny in
      the current layer}}}
\left(\frac{\left\Vert v-u\right\Vert
      ^{2}}{dist(v,u)^{2}}-1\right)(v-u)
\]


Original formula:

\[
disp(u)=\sum _{{\textrm{\tiny all vertices }v\atop \textrm{\tiny in
      the current layer}}}
\left(\frac{\left\Vert v-u\right\Vert}{dist(v,u)w}-1\right)(v-u)
\]

Fruchtermann-Reingold:
\[
disp(u)=\underbrace{\sum _{{\textrm{\tiny all vertices }v\atop
      \textrm{\tiny which are adjazent to }u}}\left(\frac{\left\Vert
    v-u\right\Vert
    ^{2}}{dist(v,u)^{2}}-1\right)(v-u)}_{\textrm{attractiv
    force}}-\underbrace{\sum _{{\textrm{\tiny all vertices   }v\atop
      \textrm{\tiny in the current
        layer}}}\left(\frac{0.05l^2}{\left\Vert u-v\right\Vert
    ^{2}}\right)(v-u)}_{\textrm{repulsiv force}}
\]

Original formula:

\[
disp(u)=\sum _{{\textrm{\tiny all vertices }v\atop
    \textrm{\tiny which are adjazent to }u}}
\left(\frac{\left\Vert
  v-u\right\Vert^{2}}{w^{2}}-1\right)(v-u)-
\sum _{{\textrm{\tiny all vertices}v\atop \textrm{\tiny in the current layer}}}
\left(\frac{0.05w^2}{\left\Vert u-v\right\Vert
  ^{2}}\right)(v-u)
\]

where $u,v$ are the position vectors of the corresponding vertices,
$disp(u)$ states the direction in which the vertex $u$ is moved
during an iteration%
\footnote{The vector $disp(u)$ is normalized before use; only the 
change of the direction compared to the previous iteration influences
the amplitude strength. %
}, $\left\Vert v-u\right\Vert $ is the current distance in the coordination
system, $dist(v,u)$ the required distance given by the weight of
the edges, $l$ is the average edge weight and $w=32$ the constant edge
weight of the original version. The scaling factor 0.05 was
taken from the original version; this value seems to provide the best
results. In the calculation of $disp(u)$  according to Kamada-Kawai,
$dist(v,u)$ is simply replaced by the given edge length instead of the 
constant. However, the calculation
of the attractive force of Fruchtermann-Reingold now considers
$dist(v,u)$ (which it did not before).
\end{itemize}

No changes were made in the calculation of the layers or of the local
temperature.


\section{Technical Details}


\subsection{Source Code Overview}

\begin{longtable}{|p{3.5cm}|p{9.5cm}|}
\hline 
File&
Description\\
\hline
\hline 
main.c&
main program - combines creation/reading of the graph and drawing\\
\hline 
gml\_{*}.c&
implementation of gml-file-reader (converter, demo, parser, scanner)\\
\hline 
GraphMLConverter.cc&
implementation of GraphML-file-reader\\
&
written in C++ because of the xerces library\\
\hline 
GraphMLException.cc,&
handles parsing errors for GraphML-files\\
handlers.cc&
\\
\hline 
GraphMLNames.cc, GDNames.cc, GD.h&
definition of the standard GraphML-tags and their extensions 
for handling the weight of an edge and the position of a vertex\\
\hline 
GraphMLDocument.cc, &
further tools for the GraphML-Converter\\
GraphML.h, &
\\
DOMAdmin.cc/.h&
\\
\hline 
cplusplusinterface.h&
declaration of functions written in C++ that are called from c\\
\hline 
bfs.c/.h, dijkstra.c/.h&
implementation of a graph breadth-first-search- and Dijkstra's algorithm\\
\hline 
Queue.c, PQueue.c, &
implementation of a queue, a prioritiy queue and a heap \\
Heap.c&
(which is used by the priority queue) needed for bfs and Dijkstra\\
\hline 
Point.c/.h&
implementation of a point-data structure\\
\hline 
Graph.c&
implementation of test graphs like trees, (un)weighted grids, cubes,
...\\
\hline 
Graph.h&
definition of the main data structure {}``graph''\\
\hline 
misf.c/.h&
creation of the misf-layers\\
\hline 
DrawGraph.c&
implementation of the vertex-positioning algorithms\\
\hline 
MesaPlot.c/.h&
implementation of plotting and user interface specific tasks\\
\hline 
graph.tcl&
tcl-user interface -- makes it easy to set the parameters for the program\\
\hline 
help.data, help.tcl&
online help for the tcl-interface\\
\hline 
graphml-structure-1.0rc.xsd&
schema definition of the GraphML-Format
(it is not the original version but a modified version that works)%
\footnote{The changes in the schema definition do not affect the format of
the XML-file.%
}\\
\hline 
Makefile&
for building the source code\\
\hline 
readme&
some information for building and executing the program\\
\hline
\end{longtable}


\subsection{Requirements}

\begin{itemize}
\item the programm was tested on Linux SuSE 8.0/7.3-Systems but should work
with other distributions as well 
\item xerces version 1.7.0 (1.6.0 might work too but you need to adjust
include paths in the source code) 
\item an OpenGL-library (e.g. MesaGlu)
\end{itemize}

\subsection{Building}

You have to set the correct paths to your installation of xerces in
the Makefile. After that type \texttt{\char`\"{}make\char`\"{}}.


\subsection{Running }

Probably you have to adjust the environment variable \texttt{LD\_LIBRARY\_PATH}
to include the xerces library.\footnote{e.g. \texttt{\char`\"{}export LD\_LIBRARY\_PATH=\char`\"{}/home/student/maagv/hiwi/xerces-c1\_6\_0-linux/lib:\$LD\_LIBRARY\_PATH\char`\"{}}} 

Then type \texttt{\char`\"{}graph.tcl\char`\"{}} or \texttt{\char`\"{}wish graph.tcl\char`\"{}}
to set the parameters with an interface or start \texttt{main} directly setting
the parameters ``by hand''. If results are written,
you can find them in the file \texttt{erg.xml}. 

You may want to use the help provided by graph.tcl


\subsection{Parsing GraphML-Files}

The validation of the xerces parser works only with an modified version
of the GraphML schema, but you can switch off validation by compiling with the
compiler option -DNO\_XERCES\_VAL. The program reads the schema style
GraphML format but it should be able to validate the former dtd style
format, too. While writing, the program tries to keep up all information
it read (if it read a GraphML-file) but e.g. the direction property
will be lost. Furthermore it may be suitable for large graphs to
work with GraphML files with no further information because all this
information will be kept in memory, so less memory will be available
for computing. Lack of memory is not caught and will propably result
in a segmentation fault.


\subsection{Specification of the allowed graphs to be parsed}

The programm accepts GraphML-graphs with the following properties: 

\begin{itemize}
\item undirected graphs 
\item directed graphs; \noun{restriction}: will be converted in undirected
graphs (for computing and writing) 
\item weighted, unweighted or partially weighted edges with a given standard
weight (If the standard weight is missing, the value 32 is used.) 
\item graphs with loops and double edges; \noun{restriction}: these will
be lost (for computing and writing) 
\item graphs with additional information; \noun{remark}: the GraphML-format
allows for every element (edge, node, the graph itself) of the graph
to have additional information. This information should be kept
when saving the graph with its new coordinates under the restrictions
noted above (but this feature was never really completely tested). 
\end{itemize}
Further restrictions: 

\begin{itemize}
\item Only the first graph in a file will be handled.
\end{itemize}
The GML-Graphs might have more properties with less restrictions,
see the documentation for original Guide.


\subsection{Specification of the extended GraphML-format}

The standard GraphML-format is not changed but extended in order to
allow the weight of edges and the coordinates of vertices to be saved.

If you want to set a weight of the edges, insert the following lines
after the GraphML tag:

\begin{lyxcode}
<key~id=\char`\"{}we\char`\"{}~for=\char`\"{}edge\char`\"{}~attr.name=\char`\"{}weight\char`\"{}~attr.type=\char`\"{}double\char`\"{}>

<default>\emph{default~weight~of~edge}</default>

</key>
\end{lyxcode}
The number enclosed by the tag \texttt{default} defines a default weight. 
Setting a default weight is optional. To set the weight of an edge, insert 
a data-tag to an edge-tag like this:

\begin{lyxcode}
<edge~source=\char`\"{}v0\char`\"{}~target=\char`\"{}v1\char`\"{}>~

<data~key=\char`\"{}we\char`\"{}>\emph{weight~of~edge}</data>~

</edge>
\end{lyxcode}
The weight of an edge can be given as an integer or a real value.

The coordinates are stored like this:

\begin{lyxcode}
<node~id=\char`\"{}v0\char`\"{}>~

<data~key=\char`\"{}x1\char`\"{}>~50.1536~</data>~

<data~key=\char`\"{}x2\char`\"{}>~168.311~</data>~

<data~key=\char`\"{}x3\char`\"{}>~-62.3714~</data>

</node>
\end{lyxcode}
If there are four or two coordinates, four or two coordinates will
be saved. If there are coordinates in the file, the following key-tags
must also be present:

\begin{lyxcode}
<key~id=\char`\"{}x1\char`\"{}~for=\char`\"{}node\char`\"{}~attr.name=\char`\"{}x1-coord\char`\"{}~attr.type=\char`\"{}double\char`\"{}/>~~

<key~id=\char`\"{}x2\char`\"{}~for=\char`\"{}node\char`\"{}~attr.name=\char`\"{}x2-coord\char`\"{}~attr.type=\char`\"{}double\char`\"{}/>~~

<key~id=\char`\"{}x3\char`\"{}~for=\char`\"{}node\char`\"{}~attr.name=\char`\"{}x3-coord\char`\"{}~attr.type=\char`\"{}double\char`\"{}/>\end{lyxcode}

\end{document}
